\section{Linux Investigations}

Linux systems are common targets in enterprise environments, especially in cloud, container, and server workloads. 
Understanding how to investigate Linux endpoints is essential for CTI analysts, threat hunters, and detection engineers. 
While SIEM telemetry provides high-level visibility, direct Linux commands offer immediate insight into processes, users, 
network activity, and system integrity. This chapter provides a focused set of commands that support real-world investigations, 
along with explanations of when and why each command is useful.

% --------------------------------------------------------------------
% SECTION 1: Core Linux Telemetry Sources
% --------------------------------------------------------------------

\subsection*{Core Linux Telemetry Sources}

Linux investigations rely on several key data sources:

\begin{itemize}
    \item \textbf{Process telemetry} — \texttt{ps}, \texttt{top}, \texttt{lsof}
    \item \textbf{Network telemetry} — \texttt{ss}, \texttt{netstat}
    \item \textbf{Authentication logs} — \texttt{/var/log/auth.log}, \texttt{/var/log/secure}
    \item \textbf{File integrity} — \texttt{find}, \texttt{stat}, \texttt{inotify}
    \item \textbf{Service and startup telemetry} — \texttt{systemctl}, \texttt{journalctl}
    \item \textbf{auditd} — kernel-level forensic logging
\end{itemize}

These sources form the foundation of Linux endpoint investigations.

% --------------------------------------------------------------------
% SECTION 2: Process Investigation
% --------------------------------------------------------------------

\subsection*{Process Investigation}

Investigations often begin with understanding what is running on a system. The \texttt{ps} and \texttt{top} utilities 
reveal active processes, their parent-child relationships, and resource usage. These commands help identify suspicious 
binaries, unexpected services, or processes running under unusual accounts.

\textbf{List all running processes}

\begin{lstlisting}[language=bash]
ps aux
\end{lstlisting}

To investigate a specific process, analysts often need to inspect its parent process, command-line arguments, and open files. 
The \texttt{lsof} command is particularly valuable for identifying network connections or files associated with a suspicious process.

\textbf{Inspect open files for a process}

\begin{lstlisting}[language=bash]
lsof -p <PID>
\end{lstlisting}

% --------------------------------------------------------------------
% SECTION 3: Network Investigation
% --------------------------------------------------------------------

\subsection*{Network Investigation}

Network visibility is essential for detecting unauthorized services, reverse shells, or unexpected outbound traffic. 
Commands such as \texttt{ss} and \texttt{netstat} reveal active connections, listening ports, and associated processes.

\textbf{List listening ports and associated processes}

\begin{lstlisting}[language=bash]
ss -tulnp
\end{lstlisting}

This helps identify suspicious listeners, unexpected services, or malware establishing command-and-control channels.

% --------------------------------------------------------------------
% SECTION 4: Authentication Investigation
% --------------------------------------------------------------------

\subsection*{Authentication Investigation}

Authentication and privilege misuse are common indicators of compromise. Linux logs authentication events in 
\texttt{/var/log/auth.log} (Debian-based systems) or \texttt{/var/log/secure} (RHEL-based systems).

\textbf{View recent authentication events}

\begin{lstlisting}[language=bash]
grep "session" /var/log/auth.log
\end{lstlisting}

This reveals logons, sudo usage, SSH sessions, and privilege escalation attempts.

% --------------------------------------------------------------------
% SECTION 5: File Integrity and Persistence
% --------------------------------------------------------------------

\subsection*{File Integrity and Persistence}

Attackers often modify system binaries, create persistence mechanisms, or drop malicious files. 
The \texttt{find} command helps locate recently modified files, which can reveal tampering or unauthorized activity.

\textbf{Find recently modified files}

\begin{lstlisting}[language=bash]
find / -type f -mtime -1 2>/dev/null
\end{lstlisting}

System services and startup mechanisms are frequent targets for persistence. The \texttt{systemctl} command provides insight 
into active services, failed services, and recently modified units.

\textbf{List active services}

\begin{lstlisting}[language=bash]
systemctl list-units --type=service
\end{lstlisting}

% --------------------------------------------------------------------
% SECTION 6: auditd and Forensic Logging
% --------------------------------------------------------------------

\subsection*{auditd and Forensic Logging}

For deeper forensic analysis, \texttt{auditd} provides detailed system call logging, including file access, privilege escalation 
attempts, and execution events. When enabled, audit logs can be queried directly.

\textbf{Search for execution events}

\begin{lstlisting}[language=bash]
ausearch -m EXECVE
\end{lstlisting}

auditd is particularly useful for detecting:

\begin{itemize}
    \item privilege escalation
    \item unauthorized file access
    \item suspicious command execution
    \item persistence mechanisms
\end{itemize}

% --------------------------------------------------------------------
% SECTION 7: SIEM Pivots (SPL + KQL)
% --------------------------------------------------------------------

\subsection*{Pivoting to SIEM Telemetry}

Linux investigations often require correlating endpoint findings with SIEM telemetry. 
For example, a suspicious process discovered with \texttt{ps aux} can be cross-referenced with Sysmon for Linux 
or audit logs ingested into Splunk or Sentinel.

\textbf{SPL Example}

\begin{lstlisting}
index=linux_logs process_name="bash"
| table _time, host, user, parent_process, command
\end{lstlisting}

\textbf{KQL Equivalent}

\begin{lstlisting}
LinuxSysmon
| where ProcessName == "bash"
| project TimeGenerated, HostName, User, ParentProcessName, CommandLine
\end{lstlisting}

These pivots allow analysts to connect endpoint-level findings with SIEM-level visibility.

% --------------------------------------------------------------------
% SECTION 8: Why Linux Investigations Matter
% --------------------------------------------------------------------

\subsection*{Why Linux Investigations Matter}

Linux systems underpin critical infrastructure, cloud workloads, and enterprise applications. 
Mastering Linux investigations provides visibility into execution, authentication, persistence, and network behavior. 
These techniques form the foundation of Linux threat hunting and incident response. Later chapters expand these fundamentals 
into full hunting playbooks and detection patterns tailored to Linux environments.
